\paragraph{Problem 2 answer}\GradescopeComment{This should be on page 6 of your PDF.}

\begin{enumerate}[label=(\alph*)]

\item  Nodes in level 0: $a$

  Nodes in level 1: \BLUE{$
    \REPLACEME
    $}

  Nodes in level 2: \BLUE{$
    \REPLACEME
    $}

  Nodes in level 3: \BLUE{$
    \REPLACEME
    $}

  Nodes in level 4: \BLUE{$
    \REPLACEME
    $}

\item

  Parent of $b$: \BLUE{$
    \REPLACEME
    $}

  Parent of $c$: \BLUE{$
    \REPLACEME
    $}

  Parent of $d$: \BLUE{$
    \REPLACEME
    $}

  Parent of $e$: \BLUE{$
    \REPLACEME
    $}

  Parent of $f$: \BLUE{$
    \REPLACEME
    $}

  Parent of $g$: \BLUE{$
    \REPLACEME
    $}

  Parent of $h$: \BLUE{$
    \REPLACEME
    $}
\end{enumerate}

\newpage

\paragraph{Problem 2(c) answer (OPTIONAL!)}\footnote
{\color{blue}
  SEE FURTHER INSTRUCTIONS IN THE COMMENTS IN PROBLEM\_2\_ANSWER.TEX!}
\GradescopeComment{This should be on page 7 (top) of your PDF.}

\begin{blue}
  
% HERE ARE TWO WAYS TO DRAW A ROOTED TREE IN LaTeX.
% UNCOMMENT ONE OF THEM AND EDIT IT TO MAKE YOUR ANSWER.

%%%%%%%%%%%%%%%%%%%%%%%%%%%%%%%%%%%%%%%%%%%%%%%%%%%%%%
%%% METHOD 1

%   \begin{tikzpicture}[  mystyle/.style ={circle, draw, minimum width = 0.6 cm}  ]

%     % draw and name the nodes 

%     \node[mystyle, label={0:a}] (R) {r};   % draw a node named R with label r

%     \node[mystyle, label={0:b}] (S) [below =of R, xshift=-50pt] {s};
%     % this line draws a node named S with center label "s" and side label "b"
%     % "0" is the rotation of the side label around the node
%     % "b" is the side label
%     % "S" is the name that can be used below to reference this node 
%     % "R" is the previously placed node that this node is placed below
%     % "-50pt" controls the horizontal placement of this node
%     % "s" is the center label (displayed within this node)
    
%     \node[mystyle, label={0:c}] (T) [below =of R, xshift= 0pt] {t};
%     \node[mystyle, label={0:d}] (U) [below =of R, xshift=50pt] {u};
%     \node[mystyle, label={0:e}] (V) [below =of S, xshift=-25pt] {v};
%     \node[mystyle, label={270:f}] (X) [below =of S, xshift=25pt] {x};

%     % draw the edges
%     \path
%     (R) edge (S)
%     (R) edge (T) 
%     (R) edge (U)
%     (S) edge (V) 
%     (S) edge (X); 
%   \end{tikzpicture}

%%%%%%%%%%%%%%%%%%%%%%%%%%%%%%%%%%%%%%%%%%%%%%%%%%%%%%
%%%   METHOD 2
  
%  % DRAW IT ON PAPER, SCAN IT TO A PDF, THEN INCLUDE THE PDF IN YOUR DOCUMENT AS FOLLOWS:

%   \includegraphics[height=3in,width=6in]{hwk_5_demo_scan}
  
%   % SEE https://texblog.org/2012/02/23/crop-figures-with-includegraphics/ IF YOU WANT TO CROP THE INCLUDED IMAGE.
  
\end{blue}

\vfill

\paragraph{Problem \arabic{problem} collaboration and citations}~\GradescopeComment{This should be on page 7 (bottom) of your PDF.}

I collaborated on this problem with the following people:
\begin{blue}
  \begin{itemize}
  \item \REPLACEME                    % replace with "none" if none
  \end{itemize}
\end{blue}

My answers use ideas from the following sources (other than the lecture notes and textbooks): 
\begin{blue}
  \begin{itemize}
  \item \REPLACEME                    % replace with "none" if none
  \end{itemize}
\end{blue}



%%% Local Variables:
%%% mode: latex
%%% TeX-master: "homework_5"
%%% End:
